\chapter{Quadratic dynamics and finite observations}
\label{ch:quadratic-foundations}
\label{chap:quadratic-foundations}

The same polynomial acts in two different roles.  With $c$ fixed, it
iterates points in the dynamical plane.  With the initial point fixed at
the critical point, it defines functions of the parameter $c$.  Keeping
these roles separate is essential when passing from Julia sets to the
Mandelbrot set.

\section{The family and its distinguished orbit}

\begin{definition}[Quadratic family]
\label{def:quadratic-family}
For $c,z\in\C$ and $n\in\N=\{0,1,\ldots\}$, set
\[
 f_c(z)=z^2+c,\qquad
 f_c^{\circ0}(z)=z,\qquad
 f_c^{\circ(n+1)}(z)=f_c(f_c^{\circ n}(z)).
\]
The critical orbit is the sequence
\[
 p_n(c)=f_c^{\circ n}(0).
\]
\end{definition}

Since $f_c'(z)=2z$, the unique critical point is $0$ and its critical
value is $c$.  In particular,
\[
 p_0(c)=0,\quad p_1(c)=c,\quad
 p_2(c)=c^2+c,\quad p_3(c)=(c^2+c)^2+c.
\]
The orbit starting at $c$ is the critical orbit with its first term
removed:
\[
 f_c^{\circ n}(c)=p_{n+1}(c).
\]
Thus bounds on either orbit differ by at most the initial value $0$.
This shift explains several appearances of the critical value, rather
than the critical point, in phase--parameter constructions.

\begin{lemma}[Dependence on the parameter]
\label{lem:orbit-polynomial}
For each $n$, $p_n$ is a polynomial in $c$, hence continuous.  For $n\ge1$
it is monic of degree $2^{n-1}$.
\end{lemma}

\begin{proof}
The assertion starts with $p_1(c)=c$.  The recurrence
$p_{n+1}=p_n^2+c$ preserves polynomial dependence.  When $n\ge1$,
the leading term of $p_n^2$ has degree $2^n>1$, so adding $c$ changes
neither its coefficient nor its degree.
\end{proof}

\paragraph{Lean correspondence.}
The review names
\leanref{Mlc/Literature/Foundations.lean}{quadraticMap} and
\leanref{Mlc/Literature/Foundations.lean}{quadraticOrbit}
are abbreviations of the imported quadratic map and orbit.  The recurrence
is \lean{MLC.Quadratic.orbit_succ}; continuity in the parameter is
\lean{Molecule.continuous_orbit}.  The degree calculation above is a
text proof; a named review lemma for that degree remains to be added.

\section{Bounded orbits and escape}

\begin{definition}[Bounded orbit]
An orbit is bounded if there is $B\in\R$ such that
\[
 |f_c^{\circ n}(z)|\le B\qquad(n\in\N).
\]
It escapes if $|f_c^{\circ n}(z)|\to+\infty$.
\end{definition}

Unboundedness of an arbitrary sequence does not imply convergence to
infinity.  For quadratic orbits the following growth estimate supplies
that implication.

\begin{lemma}[Quantitative escape estimate]
\label{lem:escape-growth}
Suppose $a=|z|>2$ and $a\ge|c|$.  Set $q=a-1>1$.  Then
\[
 |f_c^{\circ n}(z)|\ge a q^n\qquad(n\in\N).
\]
\end{lemma}

\begin{proof}
For any $w$ with $r=|w|\ge a$, the reverse triangle inequality gives
\[
 |f_c(w)|\ge r^2-|c|
       \ge r^2-r=r(r-1)\ge qr.
\]
The estimate at $n=0$ is equality.  Induction now gives the stated bound,
and in particular keeps every subsequent modulus at least $a$.
Since $q>1$, the lower bound tends to infinity.
\end{proof}

\begin{theorem}[Escape-radius criterion]
\label{thm:escape-radius}
If $|z|>\max\{2,|c|\}$, then the orbit of $z$ escapes.  More generally,
every unbounded quadratic orbit escapes.
\end{theorem}

\begin{proof}
The first assertion follows from Lemma~\ref{lem:escape-growth}.  An
unbounded orbit has a term of modulus greater than
$\max\{2,|c|\}$.  Apply the first assertion to that term and discard the
finite initial segment.
\end{proof}

The multiplicative lower bound is the important point in this proof.
Strictly increasing moduli alone would allow a finite limit and would
therefore leave a gap in the escape argument.

\begin{definition}[Filled Julia and Julia sets]
\label{def:filled-julia}
The filled Julia set is the bounded-orbit set
\[
 K_c=\{z:\exists B\in\R\ \forall n,\ |f_c^{\circ n}(z)|\le B\},
\]
and the Julia set is its topological boundary $J_c=\partial K_c$.
\end{definition}

\begin{definition}[Mandelbrot set]
\label{def:mandelbrot}
The Mandelbrot set is
\[
 \M=\{c:0\in K_c\}
    =\{c:\exists B\in\R\ \forall n,\ |p_n(c)|\le B\}.
\]
\end{definition}

\begin{theorem}[Universal critical-orbit bound]
\label{thm:critical-bound}
For every $c\in\C$,
\[
 c\in\M\quad\Longleftrightarrow\quad
 |p_n(c)|\le2\text{ for every }n.
\]
Consequently $\M\subseteq\overline B(0,2)$.
\end{theorem}

\begin{proof}
If $|c|>2$, the critical value $p_1(c)=c$ satisfies the hypotheses of
Lemma~\ref{lem:escape-growth}, including the allowed equality
$|p_1(c)|=|c|$.  Thus $c\notin\M$.
Suppose instead that $|c|\le2$ and some $p_j(c)$ has modulus greater
than $2$.  Then $|p_j(c)|>|c|$, so the orbit again escapes.
These two cases prove that a bounded critical orbit stays in the closed
disk of radius $2$.  The converse follows by choosing the bound $B=2$,
and the disk inclusion follows from $p_1(c)=c$.
\end{proof}

\paragraph{Lean correspondence.}
Boundedness is expressed exactly by
\leanref{Mlc/Literature/Foundations.lean}{filledJulia_iff_bounded_orbit}
and
\leanref{Mlc/Literature/Foundations.lean}{criticalOrbit_iff_mandelbrot}.
The universal radius statement is the imported theorem
\lean{Molecule.mandelbrot_eq_inter}; its use is visible in
\leanref{Mlc/CategoricalMandelbrot.lean}{set_subset_outerOrbitSet}.
The distinction between an escaping and a bounded critical orbit is
packaged by
\leanref{Mlc/ParameterClassification.lean}{not_mem_mandelbrot_iff_exists_escape}.

\section{Compactness and invariance in the dynamical plane}

\begin{proposition}[A closed-set description of $K_c$]
\label{prop:filled-julia-compact}
With $R_c=\max\{2,|c|\}$,
\[
 K_c=\bigcap_{n\ge0}(f_c^{\circ n})^{-1}
                       \bigl(\overline B(0,R_c)\bigr).
\]
In particular $K_c$ is compact.
\end{proposition}

\begin{proof}
A bounded orbit cannot contain a term of modulus greater than $R_c$,
by the escape criterion.  Conversely the displayed inequalities give
a common bound.  Every iterate is continuous, so the intersection is
closed.  The term $n=0$ places the intersection in
$\overline B(0,R_c)$.  A closed bounded subset of $\C\simeq\R^2$ is
compact by the Heine--Borel theorem.
\end{proof}

\begin{proposition}[Complete invariance]
\label{prop:filled-julia-invariant}
For every $c$,
\[
 f_c^{-1}(K_c)=K_c,\qquad f_c(K_c)=K_c.
\]
\end{proposition}

\begin{proof}
The orbit of $f_c(z)$ is the tail of the orbit of $z$.  Adding or removing
the single finite initial value $z$ preserves boundedness.  This proves
the inverse-image equality and the inclusion $f_c(K_c)\subseteq K_c$.
For $w\in K_c$, a solution of $z^2=w-c$ exists in $\C$.  Its image is
$w$, and the inverse-image equality places $z$ in $K_c$, proving the
reverse inclusion.
\end{proof}

These are statements about a fixed dynamical system.  Connectedness of
$K_c$ and variation of $K_c$ with $c$ require additional arguments.
The classical connectedness criterion
\[
 K_c\text{ connected}\quad\Longleftrightarrow\quad c\in\M
\]
belongs to the external-coordinate theory of
Douady--Hubbard~\cite{DH}.  It is distinct from local connectivity of
either $J_c$ or $\M$.

\section{Periodic orbits, multipliers, and the first components}

\begin{definition}[Periodic point and multiplier]
A point $z$ is periodic with return period $p\ge1$ if
$f_c^{\circ p}(z)=z$.  Its prime period is the least such $p$.
The multiplier for a return period is
\[
 \lambda_{c,p}(z)=(f_c^{\circ p})'(z).
\]
The cycle is attracting, repelling, or neutral according as the
modulus of its prime-period multiplier is less than, greater than,
or equal to $1$.  It is superattracting when that multiplier is zero.
\end{definition}

\begin{proposition}[Multiplier product]
For any $z$ and $p\ge1$,
\[
 (f_c^{\circ p})'(z)
   =\prod_{j=0}^{p-1}2f_c^{\circ j}(z).
\]
For a periodic point, changing the starting point on its cycle
preserves this product.  Passing from prime period $p$ to return
period $mp$ raises the multiplier to its $m$-th power.
\end{proposition}

\begin{proof}
Apply the chain rule repeatedly, using $f_c'(w)=2w$.
On a periodic orbit a change of starting point cyclically permutes
the factors.  A return through the cycle $m$ times repeats the
same product $m$ times.
\end{proof}

\begin{definition}[Parabolic, Siegel, and Cremer cycles]
A neutral cycle is parabolic when its multiplier is a root of
unity.  At an irrationally neutral cycle, the multiplier has the
form $e^{2\pi i\theta}$ with $\theta\notin\Q$.
Such a cycle is Siegel when its return map is locally analytically
conjugate to $w\mapsto e^{2\pi i\theta}w$, and Cremer when no such
local conjugacy exists.  A local analytic conjugacy is a conformal
coordinate near the periodic point that intertwines the two maps.
\end{definition}

\begin{definition}[Hyperbolic and Misiurewicz parameters]
For $c\in\M$, call $c$ hyperbolic if the critical orbit is attracted
to an attracting periodic cycle $\mathcal O$, meaning
$\min_{w\in\mathcal O}|p_n(c)-w|\to0$.
Call it Misiurewicz if the
critical point is strictly preperiodic to a repelling cycle.
Here strictly preperiodic means that the orbit eventually becomes
periodic but its starting point is not periodic.
In this quadratic connectedness-locus setting, the subhyperbolic
cases consist of the hyperbolic and Misiurewicz cases.
\end{definition}

\paragraph{Lean correspondence.}
The return-period predicate and derivative are
\leanref{Mlc/ParameterClassification.lean}{PeriodicPoint} and
\leanref{Mlc/ParameterClassification.lean}{multiplier}.
The structures and predicates \lean{LocalLinearization},
\lean{HasParabolicCycle}, \lean{HasSiegelCycle}, and
\lean{HasCremerCycle} express the indicated local conditions.
The chain-rule product is formal as
\leanref{Mlc/Quadratic/Complex/FirstComponents.lean}{deriv_fc_iterate};
the invariance of this product under a change of starting point on a
cycle remains a text argument.
The theorem \lean{classification_complete} is an exhaustive
priority partition with an explicit residual case; it neither
proves MLC on each part nor classifies all renormalization types.

\begin{proposition}[The attracting fixed-point component]
\label{prop:main-cardioid}
The parameters with an attracting fixed point form
\[
 \mathcal H_0=
 \left\{\frac{\lambda}{2}-\frac{\lambda^2}{4}:
                         |\lambda|<1\right\}.
\]
The boundary parametrization is
$c(\theta)=e^{2\pi i\theta}/2-e^{4\pi i\theta}/4$.
Its cusp $c(0)=1/4$ has fixed-point multiplier $1$.
The parametrization is injective on the closed unit disk, an attracting
fixed point of $f_c$ is unique when it exists, the set $\mathcal H_0$ is
open, and its frontier is the image of the unit circle.
\end{proposition}

\begin{proof}
If $\alpha$ is a fixed point, then
$c=\alpha-\alpha^2$ and its multiplier is $\lambda=2\alpha$.
Conversely, the displayed parameter has fixed point $\lambda/2$
with multiplier $\lambda$.  The map
$\lambda\mapsto\lambda/2-\lambda^2/4$ is injective on the closed
unit disk: for $\lambda\ne\mu$, equality forces $\lambda+\mu=2$, and
two points of the closed disk with sum $2$ both equal $1$.
If $\alpha\ne\beta$ are fixed points, then $\alpha+\beta=1$.  If both were
attracting, then $2=|2\alpha+2\beta|\le|2\alpha|+|2\beta|<2$, a contradiction.
The set $\mathcal H_0$ is the image of the open disk under a nonconstant
polynomial, so it is open by the open-mapping theorem.  The image of the
closed disk is compact, and the parametrization is injective there, so the
closure of $\mathcal H_0$ is that image and its frontier is the image of the
unit circle $\{e^{2\pi i\theta}\}$.  This gives the displayed boundary curve.
\end{proof}

{\rightskip=0pt plus 1fil
\paragraph{Lean correspondence.}
The attracting-fixed-point predicate is
\leanref{Mlc/Quadratic/Complex/FirstComponents.lean}{HasAttractingFixedPoint}.
The description of $\mathcal H_0$ is
\leanref{Mlc/Quadratic/Complex/FirstComponents.lean}{hasAttractingFixedPoint_iff_mem_image}.
Injectivity on the closed disk is
\leanref{Mlc/Quadratic/Complex/FirstComponents.lean}{cardioidMap_injOn_closedBall},
uniqueness is
\leanref{Mlc/Quadratic/Complex/FirstComponents.lean}{attracting_fixed_point_unique},
openness is
\leanref{Mlc/Quadratic/Complex/FirstComponents.lean}{isOpen_image_ball_cardioid},
and the frontier identity is
\leanref{Mlc/Quadratic/Complex/FirstComponents.lean}{frontier_hasAttractingFixedPoint}.\par}

\begin{proposition}[Multiplier map of the main cardioid]
\label{prop:multiplier-cardioid}
Let $\sqrt{\cdot}$ denote the principal square root, defined on $\C\setminus(-\infty,0]$ and holomorphic there.  For $c\in\mathcal H_0$ the number $1-4c$ is not in $(-\infty,0]$, and the attracting fixed point of $f_c$ has multiplier
\[
 \lambda(c)=1-\sqrt{1-4c}.
\]
The map $\lambda\colon\mathcal H_0\to\mathbb D$ is a biholomorphism, with inverse $t\mapsto t/2-t^2/4$.  This is the period-one case of the multiplier theorem in Lyubich's notes.
\end{proposition}

\begin{proof}
Let $c\in\mathcal H_0$ with attracting fixed point of multiplier $t$, so that $|t|<1$ and $t/2-t^2/4=c$.  Then $(1-t)^2=1-4c$.  Write $s=\sqrt{1-4c}$, so $(1-t-s)(1-t+s)=0$.  If $1-t=-s$, then $t=1+s$ has real part at least $1$, since $\operatorname{Re}s\ge0$; this contradicts $|t|<1$.  Hence $t=1-s=\lambda(c)$.

The point $1-4c$ lies in the slit plane: if $c$ is real then $c\in(-\tfrac34,\tfrac14)$ by Corollary~\ref{cor:real-components}, and otherwise $\operatorname{Im}(1-4c)\neq0$.  So $\lambda$ is holomorphic on $\mathcal H_0$.  For $|t|<1$ the parameter $c=t/2-t^2/4$ lies in $\mathcal H_0$, and the argument above gives $\lambda(c)=t$; this is the identity $\lambda\circ(t\mapsto t/2-t^2/4)=\mathrm{id}$.  Conversely, for $c\in\mathcal H_0$ the identity $(1-\lambda(c))^2=1-4c$ gives $\lambda(c)/2-\lambda(c)^2/4=c$; this is the identity $(t\mapsto t/2-t^2/4)\circ\lambda=\mathrm{id}$.  Both maps are holomorphic, so $\lambda$ is a biholomorphism.
\end{proof}

\paragraph{Lean correspondence.}
The statement is
\leanref{Mlc/Quadratic/MultiplierMap.lean}{multiplier_map_biholomorphic}, a conjunction of the mapping properties, the two inverse identities, and differentiability.  The square-root identity is \leanref{Mlc/Quadratic/MultiplierMap.lean}{mainSqrt_sq}, the sign of the real part is \leanref{Mlc/Quadratic/MultiplierMap.lean}{re_mainSqrt_nonneg}, and the identification of the multiplier is \leanref{Mlc/Quadratic/MultiplierMap.lean}{eq_mainMultiplier_of_cardioidMap}.


\begin{proposition}[The attracting period-two component]
\label{prop:period-two-component}
The parameters with an attracting cycle of prime period two form
the disk
\[
 \mathcal H_2=\{c:|c+1|<1/4\}.
\]
On this disk the cycle multiplier is $4(c+1)$, and the frontier of
$\mathcal H_2$ is the circle $|c+1|=1/4$.
\end{proposition}

\begin{proof}
The factorization
\[
 f_c^{\circ2}(z)-z
 =(z^2-z+c)(z^2+z+c+1)
\]
separates fixed points from period-two candidates.  If $z$ has prime
period two, then $f_c(z)\ne z$ forces $z^2+z+c+1=0$.  The identity
\[
 z\,f_c(z)-(c+1)=(z-1)(z^2+z+c+1)
\]
gives $z\,f_c(z)=c+1$, and the chain rule gives the cycle multiplier
$(f_c^{\circ2})'(z)=4z\,f_c(z)=4(c+1)$.
Conversely, let $|c+1|<1/4$, and let $z$ be a root of $z^2+z+c+1$.  If
$z$ were fixed, subtracting $z^2+c=z$ from the root equation would give
$z=-1/2$, hence $c=-3/4$, contradicting $|c+1|<1/4$.  So $z$ has prime
period two, and it is attracting because $|4(c+1)|<1$.
\end{proof}

{\rightskip=0pt plus 1fil
\paragraph{Lean correspondence.}
The period-two predicate is
\leanref{Mlc/Quadratic/Complex/FirstComponents.lean}{HasAttractingPeriodTwoCycle}.
The factorization is
\leanref{Mlc/Quadratic/Complex/FirstComponents.lean}{period_two_factorization},
the identity $z\,f_c(z)=c+1$ is
\leanref{Mlc/Quadratic/Complex/FirstComponents.lean}{period_two_mul},
the multiplier formula is
\leanref{Mlc/Quadratic/Complex/FirstComponents.lean}{period_two_multiplier},
the disk description is
\leanref{Mlc/Quadratic/Complex/FirstComponents.lean}{hasAttractingPeriodTwoCycle_iff},
and the frontier is
\leanref{Mlc/Quadratic/Complex/FirstComponents.lean}{frontier_hasAttractingPeriodTwoCycle}.\par}

These formulas identify two open covered regions and their classical
boundaries; they explain the green contours in Figure~\ref{fig:mlc-status}.
The formal sets are defined by attracting fixed points and attracting
period-two cycles.  Identifying them with components of the interior of
$\M$ needs the attracting-basin argument, which is not formalized here.
They provide no classification of the rest of $\partial\M$.

\begin{proposition}[Disjointness and the root of the bulb]
\label{prop:cardioid-bulb-contact}
The sets $\mathcal H_0$ and $\mathcal H_2$ are disjoint.  Their closures meet
in exactly one point, the parameter $c=-3/4$.
\end{proposition}

\begin{proof}
Let $c=\lambda/2-\lambda^2/4$ with $|\lambda|<1$.  Then
\[
 4(c+1)=5-(\lambda-1)^2.
\]
Since $|\lambda-1|\le|\lambda|+1<2$, the reverse triangle inequality gives
$|4(c+1)|\ge5-|\lambda-1|^2>1$.  Hence $|c+1|>1/4$, and $c\notin\mathcal H_2$.

Now let $c$ lie in both closures.  Then $c=\lambda/2-\lambda^2/4$ with
$|\lambda|\le1$, and $|c+1|\le1/4$ gives $|5-(\lambda-1)^2|\le1$.  Hence
$|(\lambda-1)^2|\ge4$, that is, $|\lambda-1|\ge2$.  Write $\lambda=x+iy$.  The
two inequalities $x^2+y^2\le1$ and $(x-1)^2+y^2\ge4$ give $x\le-1$, while
$x^2\le1$ gives $x\ge-1$.  So $x=-1$ and $y=0$, whence $c=-3/4$.
Conversely, $\lambda=-1$ is in the closed disk and gives $c=-3/4$, and
$|-3/4+1|=1/4$ shows that $-3/4$ lies on the bulb boundary.
\end{proof}

{\rightskip=0pt plus 1fil
\paragraph{Lean correspondence.}
Disjointness is
\leanref{Mlc/Quadratic/Complex/FirstComponents.lean}{disjoint_cardioid_period_two}.
The common closure point is
\leanref{Mlc/Quadratic/Complex/FirstComponents.lean}{closure_cardioid_inter_closure_period_two},
which uses the closure identity
\leanref{Mlc/Quadratic/Complex/FirstComponents.lean}{closure_image_ball_cardioid}.\par}

\section{Finite observations in parameter space}

\begin{definition}[Finite outer approximation]
\label{def:finite-outer}
For $N\in\N$, set
\[
 O_N=\{c:|c|\le2,\quad |p_j(c)|\le2\text{ for }0\le j\le N\}.
\]
\end{definition}

The stages $O_0$ and $O_1$ both equal $\overline B(0,2)$; the first
additional orbit constraint appears at $O_2$:
\[
 O_2=\{c:|c|\le2,\ |c^2+c|\le2\}.
\]
Writing $c=x+iy$, the squared norm constraints are polynomial
inequalities in $x,y$.  Thus each finite stage is a real semialgebraic
set.  Its finite description should be distinguished from an algorithm
for classifying all its connected components at a useful complexity.

\begin{theorem}[Compact nested outer approximation]
\label{thm:finite-outer}
For every $N$,
\[
 O_N\text{ is compact},\qquad
 O_{N+1}\subseteq O_N,\qquad
 \M\subseteq O_N,\qquad
 \bigcap_N O_N=\M.
\]
\end{theorem}

\begin{proof}
Continuity of $p_j$ makes each norm inequality closed.  Their finite
intersection with $\overline B(0,2)$ is closed and bounded, hence
compact.  Increasing $N$ adds constraints and proves nesting.
The universal critical-orbit bound proves $\M\subseteq O_N$.
Finally, if $c\in O_N$ for every $N$, choose $N=j$ to obtain
$|p_j(c)|\le2$ for every $j$; Theorem~\ref{thm:critical-bound} then
gives $c\in\M$.
\end{proof}

\paragraph{Lean correspondence.}
The definition is
\leanref{Mlc/CategoricalMandelbrot.lean}{outerOrbitSet}.
The four conclusions are
\lean{isCompact_outerOrbitSet},
\lean{outerOrbitSet_antitone},
\lean{set_subset_outerOrbitSet}, and
\lean{iInter_outerOrbitSet_eq_set} in
\lean{MLC.Categorical.Mandelbrot}.
The last equality also proves closedness and compactness of $\M$,
without using MLC.

\begin{proposition}[First-escape witnesses]
\label{prop:first-escape}
If $c\notin\M$, there is a least $n$ such that $|p_n(c)|>2$.
Every such witness excludes $c$ from all outer stages $O_N$ with $N\ge n$.
\end{proposition}

\begin{proof}
Theorem~\ref{thm:critical-bound} supplies at least one violating index.
The well-ordering of $\N$ supplies the least one.  A stage with $N\ge n$
contains the violated $n$-th constraint.
\end{proof}

The least index is a mathematical existence statement.  In an exact
computation, certifying the strict inequality additionally requires
representations and error bounds for the parameter and its orbit.
The Lean predicates \lean{escapeAt} and \lean{firstEscapeAt}, and the
theorem
\leanref{Mlc/ParameterClassification.lean}{exists_firstEscape_of_not_mem},
keep the witness separate from any claim of decidability.

\section{What convergence of outer approximations gives}

\begin{lemma}[Eventual containment]
\label{lem:outer-eventual-containment}
If an open set $U\subseteq\C$ contains $\M$, then $O_N\subseteq U$ for
some $N$.
\end{lemma}

\begin{proof}
Otherwise the compact sets $F_N=O_N\setminus U$ would be nonempty and
nested.  They lie in the compact set $O_0$, so their intersection would
be nonempty.  Yet
\[
 \bigcap_N F_N=\left(\bigcap_N O_N\right)\setminus U
              =\M\setminus U=\varnothing,
\]
a contradiction.
\end{proof}

For a nonempty compact set $A$, write
$\dist(z,A)=\inf_{a\in A}|z-a|$.  For two such sets define the Hausdorff
distance by
\[
 d_{\mathrm H}(A,B)=
 \max\left\{\sup_{a\in A}\dist(a,B),\
             \sup_{b\in B}\dist(b,A)\right\}.
\]
The set $\M$ is nonempty because $0\in\M$.  Applying
Lemma~\ref{lem:outer-eventual-containment} to its open
$\varepsilon$-neighbourhood gives
\[
 d_{\mathrm H}(O_N,\M)\longrightarrow0.
\]
This compactness proof supplies no computable rate.  It also supplies
no bound on the components of $O_N\cap B(c,r)$; that is the additional
information needed for local connectivity.

\paragraph{Formal boundary.}
The exact limit and finite-stage compactness are kernel-checked as above.
The general Hausdorff-convergence argument is included here as a text
proof, rather than attributed to a declaration that merely states the
set intersection.  Chapter~\ref{chap:topological-foundations} isolates
the separate, formally proved component-buffer implication.

\section{An exhaustion by bounds is not a finite certificate}

The categorical approximation module also defines
\[
 W_N=\{c:\forall n,\ |p_n(c)|\le N\}.
\]
The index bounds the orbit values; it does not bound the number of
observations.  Consequently every $W_N$ still involves infinitely many
conditions.

\begin{proposition}[Saturation of the orbit-bound exhaustion]
\label{prop:inner-saturation}
The sets $W_N$ are increasing, $\bigcup_N W_N=\M$, and
\[
 W_N=\M\qquad\text{whenever }N\ge2.
\]
\end{proposition}

\begin{proof}
Increasing the allowed bound preserves all inequalities.  Every bounded
orbit admits an integer upper bound, proving the union equality.
Each $W_N$ is contained in $\M$ by definition.  Conversely, the
critical-orbit bound $2$ places every $c\in\M$ in every $W_N$ with
$N\ge2$.
\end{proof}

The saturation assertion is
\leanref{Mlc/CategoricalMandelbrot.lean}{innerOrbitSet_eq_set_of_two_le};
the exhaustion assertion is \lean{iUnion_innerOrbitSet_eq_set}.
Thus this two-sided presentation is an exact logical approximation
system.  Computable inner approximations require different objects,
such as the trapping certificates in
\leanref{Mlc/CertifiedTrappingRegions.lean}{Mlc/CertifiedTrappingRegions.lean}
and
\leanref{Mlc/FiniteTrappingRegions.lean}{Mlc/FiniteTrappingRegions.lean}.
Those certificates prove forward invariance of a bounded finite union of
rectangles, uniformly over a parameter box.  Iterating the invariant-set
inclusion then proves boundedness of every critical orbit in that box.
The finite certificate is the invariant-set calculation, rather than a
finite prefix of the orbit.

\section{Examples and the meaning of a render}

Some exact critical orbits are useful reference points:
\[
\begin{array}{c|l}
c & p_0(c),p_1(c),p_2(c),\ldots\\
\hline
0 & 0,0,0,\ldots\\
-1 & 0,-1,0,-1,\ldots\\
-2 & 0,-2,2,2,\ldots\\
i & 0,i,-1+i,-i,-1+i,-i,\ldots\\
1 & 0,1,2,5,\ldots
\end{array}
\]
The first four orbits are bounded.  For $c=1$, the value $5$ at time $3$
certifies escape, although the first three values have modulus at most
$2$.  Finite survival is therefore compatible with eventual escape.

An ideal exact-arithmetic render with cutoff $N$ displays a finite-stage
survival set.  A floating-point render adds numerical errors to that
finite-time limitation.  Neither version identifies the set of parameters
not yet covered by published local-connectivity theorems.  In particular,
the colour of a pixel certifies neither its infinite orbit nor its
renormalization type.  The coverage illustrations in this book use
separate symbols and captions for theorem classes and numerical geometry.
